\subsection{Variation independence}

\cite{Daw01b} defines variation independence in the following way.
For an underlying outcome space $\Tcal$ (just a set) and functions (``deterministic variables'') 
$X : \Tcal \to \C{X}$, $Y : \Tcal \to \C{Y}$ and $Z : \Tcal \to \C{Z}$ (with $\C{X},\C{Y},\C{Z}$ arbitrary spaces), define
the range of $Y$ as $R(Y) := \{Y(t) : t \in \Tcal\} \subseteq \C{Y}$ and the conditional range of $X$ given $Y=y \in R(Y)$ as
$R(X|Y=y) := \{X(t) : t \in \Tcal, Y(t) = y\}$. One can then define the notion of variation independence as follows:
\begin{align*}
  X \Indep^v_\Tcal Y | Z 
  &\iff \forall (y,z) \in R(Y,Z): R(X | y,z) = R(X|z) \\
  &\iff \forall (y,z) \in R(Y,Z), (y',z) \in R(Y,Z): R(X|Y=y,Z=z) = R(X|Y=y',Z=z) \\
  &\iff \forall z\in R(Z) : R(X,Y|z) = R(X|z) \times R(Y|z).
\end{align*}
This is easily shown to be a separoid (with ``being a function of'' as the ordering).
This allows to catch the notion that a function does not depend on some input:
If $X_i : \Tcal \to \C{X}_i$ and $f(X_1,X_2,X_3,\dots,X_n)$ is a function then
$f(X_1,\dots,X_n) \Indep^v_\Tcal X_j \given X_1,\dots,X_{j-1},X_{j+1},X_n$ means
that $f$ does not actually depend on $X_j$. Note that this is a conditional independence
statement, a marginal independence statement is much weaker. 

We can also consider an almost-sure version of variation independence, as follows.
Replace the outcome space $\Tcal$ by a measurable space $(\Omega,\Sigma)$ and a set of nullsets $\Ncal_\Omega$.
For measurable functions  
$X : \Tcal \to \C{X}$, $Y : \Tcal \to \C{Y}$ and $Z : \Tcal \to \C{Z}$ (with $\C{X},\C{Y},\C{Z}$ arbitrary measurable spaces), define
the range of $Y$ as $R(Y) := \{Y(t) : t \in \Tcal\} \subseteq \C{Y}$ and the conditional range of $X$ given $Y=y \in R(Y)$ as
$R(X|Y=y) := \{X(t) : t \in \Tcal, Y(t) = y\}$. One can then define the notion of variation independence as follows:

\Joris{Perhaps define a version of variation independence by $R(A,B|C(x)) = R(A|C(x)) \times R(B|C(x))\ M\as$?
Can we e.g.\ show that
$R(A|B(x),C(x)) = R(A|C(x))\ M\as$ implies $R(A,B|C(x)) = R(A|C(x)) \times R(B|C(x))\ M\as$?
Here $R(A|C(x)) := \{A(x') | C(x') = C(x)\}$ for $A,B,C$ functions on $\CX$.
Let $N$ be an $M$-null set such that $R(A|B(x),C(x)) = R(A|C(x))$ for all $x \in N^c := \CX \sm N$.
By definition:
$R(A,B|C(x)) = \{(A(x'),B(x')) | C(x') = C(x)\}$, and
$R(A|C(x)) = \{A(x') | C(x') = C(x)\}$, and $R(B|C(x)) = \{B(x') | C(x') = C(x)\}$.
Let $x \in N^c$.
Suppose $x'$ is such that $C(x') = C(x)$, then $A(x') \in R(A|C(x))$ and $B(x') \in R(B|C(x))$ and $(A(x'),B(x')) \in R(A,B|C(x))$.
It is clear that $R(A,B|C(x)) \ins R(A|C(x)) \times R(B|C(x))$.
For an element of $R(A|C(x)) \times R(B|C(x))$, there exist $x',x''$ such that $C(x')=C(x'')=C(x)$ and the element is of the form $(A(x'),B(x''))$.
Then in any case $B(x'') \in R(B|C(x'')) = R(B|C(x))$.
If $R(A|B(x''),C(x'')) = R(A|C(x''))$ then we could conclude that also $A(x') \in R(A|B(x''),C(x''))$ and hence that there is an $x'''$ with $C(x''')=C(x)$, $B(x''')=B(x'')$ and $A(x''') = A(x')$.
Then $(A(x'''),B(x''')) \in R(A|B(x''),C(x''))$.
This works with the forall-quantifiers.
Does it work $M\as$?
We need that $x'' \in N^c$ (to ensure that $R(A|B(x''),C(x'')) = R(A|C(x''))$).
So we need that for any element $b$ of $R(B|C(x))$, we can pick $x'' \in N^c$ with $C(x'')=C(x)$ and $B(x'') = b$.
It is not immediately obvious whether that holds!?!?!
}

\subsection{Variation Independence Markov Property}

Can we prove a variation-independence Markov property for SCMs?

Let us take a slightly extended SCM, where $\C{X}_J$ is not necessarily a Cartesian product,
and we have a partition $H = \{H_1,\dots,H_h\}$ of $J$ such that $X_{H_i} \Indep X_{J\sm H_i}$. 
We add edges in $G(M)$ between all input nodes within each $H_i$, for all $H_i$ in $H$.

Then, the $id$-separation Markov property holds for $X_J$ with respect to $G(M)_J$
(or: for $R(X_J)$ with respect to $G(M_J)$).

So the question is if it continues to hold further downstream.
For that, we can use the original proof by Lauritzen.

Let us start simple and derive a Markov property for an SCM without input nodes.
(Note: if we remove the exogenous distribution, then there is no longer any difference between the exogeonous input nodes and the exogenous random nodes, so this is actually the general case).
Hm, perhaps we should assume an SCM without exogenous random nodes instead.
Philosophically that is cleaner.

\begin{Thm}[Variation-Independence Global Markov property for Acyclic SCMs]\label{thm:vi_gmp_acyclic_scms}
  Let $M = \lt V,J,\CX,f \rt$ be a acyclic SCM without exogenous random nodes (with trivial $W=\emptyset$ and $P = \ind_\star$ that we suppressed in the notation).
  Let $G^+(M)$ denote its graph.
  $R_M(X_V \given \doit(X_J))$, the joint range of $X_V$ conditional on the input $X_J$, satisfies the directed global Markov property.
  In other words, for any subsets $A,B,C \subseteq V \cup J$ (not necessarily disjoint), we have that:
  $$A \idPerp_{G^+(M)} B \given C \implies X_A \Indep_{R_M(X_V \given \doit(X_J))} X_B \given X_C$$
\end{Thm}
\begin{proof}
  We write $G := G^+(M)$ and $R := R_M(X_V \given \doit(X_J))$.
  Note that:
  \[A \idPerp_G  B \given C,\] 
  implies that:
  \[ A \cap B \ins C.  \]
  Otherwise a trivial walk from $A$ to $B$ would be $C$-open. 
  This shows that $A\sm C$, $B \sm C$ and $C$ are pairwise disjoint.

  Induction start: $V=\emptyset$. 
  This means that $A,B,C \ins J$.
  The observation $A \cap B \ins C$ implies that for all $c \in R(X_C):$
  $$ R(X_A,X_B \given X_C=c) = \{(c_A,c_B)\} = \{c_a\} \times \{c_b\} = R(X_A \given X_C=c) \times R(X_B \given X_C=c).$$
  Hence $X_A \Indep_R X_B \given X_C$.

  We now assume that the claim holds for $\#V < n$ and arbitrary $J$.
  We consider the case $\#V=n>0$.
  Let $v$ be a node in $V$ without children (which always exists because we assumed $M$ to be acyclic).
  Write $D := \Pa^G(v)$.
  Its structural equation implies that $x_v = f_v(x_D)$.
  Hence,
  $$X_v \Indep_R X_{V \cup J \sm \{v\}} \given X_D.$$

  For the following we first note that if we have $A', B', C' \subseteq V \cup J$ with $v \notin A' \cup B' \cup C'$, 
  the induction hypothesis already implies that 
  \[A' \idPerp_{G} B' \given C' \implies X_{A'} \Indep_{R} X_{B'} \given X_{C'}.\]
  Indeed: 
    \begin{align*}
      A' \idPerp_{G} B' \given C'
      & \xRightarrow{\text{marginalization}} &  A' \idPerp_{G^{\sm v}}  B' \given C'\\
      & \xRightarrow{\text{induction}} & X_{A'} \Indep_{R_{M_{\sm \{v\}}}} X_{B'} \given X_{C'}\\
      & \xRightarrow{\text{marginalization}} & X_{A'} \Indep_{R} X_{B'} \given X_{C'}.
    \end{align*}

  In the following we will distinguish between 4 cases: 
  $v \in A \sm C$, $v \in B \sm C$, $v \in C$, $v \notin A \cup B \cup C$.
  Since $v \notin J$, this covers all possible cases.

  \begin{enumerate}[label=(\roman*)]
    \item $v \in A\sm C$. Then we can write:
    \begin{align*}
      A &= A' \dcup (A \cap C) \dcup \{v\}, \\
      B &= B' \dcup (B \cap C),
    \end{align*}
    with some disjoint $A' \ins A \sm C$ and $B' \ins B \sm C$.
    We then have the implications:
    \begin{align*}
        A \idPerp_G  B \given C & \xRightarrow{\text{decomposition}}  & A' \idPerp_G B' \given C \\
        & \xRightarrow{\text{induction}} & X_{A'} \Indep_{R}  X_{B'} \given X_C. \tag{\#1}
    \end{align*}
    On the other hand we have:
    \begin{align*}
     A \idPerp_G  B \given C & \xRightarrow{\text{decomposition}}  & A \idPerp_G   B' \given C   \\
      & \xRightarrow{\text{weak union}}  & \{v\} \idPerp_G  B' \given A' \dcup C    \\
     & \xRightarrow{(*), \text{ see below}}  & D  \idPerp_G   B' \given A' \dcup C \\ 
      & \xRightarrow{\text{induction}}  &   X_D  \Indep_{R}  X_{B'} \given X_{A'},X_C \tag{\#2}
    \end{align*}
    (*) holds since every $(A' \dcup C)$-open walk $ w \sus \cdots$ from a  $w \in D=\Pa^G(v)$ to $B'$ 
    extends to an $(A' \dcup C)$-open walk from $v$ to $B'$ via  $v \hut w \sus \cdots$, as $w$ stays a 
    non-collider in the extended walk (not in $A' \dcup C$) and $v \notin A' \dcup C$.

    The local Markov property for $v$ now implies:
    \begin{align*}
    &&X_v \Indep_{R}  X_{(V \cup J) \sm \{v\}} \given X_D\\ 
    &\xRightarrow{\text{decomposition}} & 
    X_v \Indep_{R}  X_{A'}, X_{B'}, X_C \given X_D \\
    & \xRightarrow{\text{weak union}}  & 
    X_v \Indep_{R} X_{B'} \given X_{A'}, X_C, X_D \\ 
    & \xRightarrow{\text{contraction}\eqref{eqn:bla}}  & X_v, X_D  \Indep_{R}  X_{B'} \given X_{A'}, X_C \\
    & \xRightarrow{\text{decomposition}}  & X_v \Indep_{R}  X_{B'} \given X_{A'}, X_C\\
    & \xRightarrow{\text{contraction} \eqref{eqn:bli}}  &  X_{A'},X_v \Indep_{R}  X_{B'} \given X_C\\
    & \xRightarrow{\text{redundancy}}  & X_{A'}, X_v \Indep_{R}  X_B \given X_C.
     \tag{\#3}
    \end{align*}
    From redundancy we get the implications:
    \begin{align*}
      &&X_{A\sm C} \Indep_{R} X_B \given X_{A'},X_v,X_C\\
    & \xRightarrow{\text{contraction} \eqref{eqn:bloop}}  &  
      X_{A'},X_v,X_{A\sm C} \Indep_{R}  X_B \given X_C\\
    & \xRightarrow{\text{decomposition}}  &    
    X_A \Indep_{R}  X_B \given X_C. 
    \end{align*}
    This shows the claim in the first case.

  \item $v \in B\sm C$. This follows from the first case by symmetry.

  \item $v \in C$.
     Then we can write:
    \begin{align*}
        A &= A' \dcup (A \cap C), \\
        %J \cup B &= B' \dcup ((J \cup B) \cap C),\\
        B &= B' \dcup (B \cap C),\\
        C &= C' \dcup \{v\},
    \end{align*}
    with some pairwise disjoint $A' \ins A \sm C$, 
    %$B' \ins (J \cup B) \sm C$, 
    $B' \ins B \sm C$ and
    $C' \ins C$.

    \Joris{The following is identical to the part in Patrick's proof and is just a property of (ordinary) $id$-separation.
    So we should turn it into a lemma.}
    {\color{gray}
    We now claim that:
    \[A' \idPerp_G  B' \given C'\dcup\{v\}\]
    implies that one of the following statements holds:
    \[  A' \dcup \{v\} \idPerp_G  B' \given C' \qquad \lor \qquad A' \idPerp_G B' \dcup \{v\} \given C'. \]
    Assume the contrary:
    \[  A' \dcup \{v\} \nidPerp_G  B' \given C' \qquad \land \qquad A' \nidPerp_G B' \dcup \{v\} \given C'. \]
    So there exist shortest $C'$-open walks $\pi_1$ and $\pi_2$ in $G$ such that all colliders are in $C'$:
    \[ \pi_1:\quad  A' \cup \{v\}\ni u_0 \sus \cdots \sus u_k \in B',  \]
    and:
    \[ \pi_2:\quad  A' \ni w_0 \sus \cdots \sus w_m \in (B' \dcup \{v\}).  \]
    So all non-colliders of $\pi_1$ and $\pi_2$ are outside of $C'$.
    Since we consider shortest walks and $v \notin C'$ at most an end node of $\pi_1$ and $\pi_2$ could be equal to $v$.
    Otherwise one could shorten the walk.\\
    Then note that $v \notin A'$ and $v \notin B'$, thus: $u_k \neq v$ and $w_0 \neq v$.\\
    If now $\pi_i$ does not contain $v$ as an (end) node, then $\pi_i$ would be $(C'\dcup \{v\})$-open,
    which is a contradiction to the assumption:
    \[A' \idPerp_G  B' \given C'\dcup\{v\}.\]
    So we can assume that the other end nodes equal $v$, i.e.: $u_0=v$ and $w_m =v$.\\
    Furthermore, both $\pi_1$ and $\pi_2$ are non-trivial walks, since $u_0 \neq u_k$ and $w_0 \neq w_m$.
    %Otherwise $v=u_k \in B'$ or $v=w_0 \in A'$, both of which are not true.
    %So $k,m \ge 1$.\\
    Since $v$ is childless and $k,m \ge 1$ we have that the $\pi_i$ are of the forms:
    \[ \pi_1:\quad   v \hut u_1 \sus \cdots \sus u_k,  \]
    and:
    \[ \pi_2:\quad  w_0 \sus \cdots \sus w_{m-1} \tuh v,  \]
    with $u_1, w_{m-1} \in D=\Pa^G(v)$. Then the following walk:
    \[ A' \ni  w_0 \sus \cdots \sus w_{m-1} \tuh v \hut u_1 \sus \cdots \sus u_k \in B',  \]
    is a $(C' \dcup \{v\})$-open walk from $A'$ to $B'$, in contradiction to:
    \[A' \idPerp_G B' \given C'\dcup\{v\}.\]
    So the claim:
    \[  A' \dcup \{v\} \idPerp_G  B' \given C' \qquad \lor \qquad A' \idPerp_G B' \dcup \{v\} \given C', \]
    must be true.
    }
    
    Therefore we get the implications:
    \begin{align*}
    A  \idPerp_G  B \given C 
      & \xRightarrow{\text{decomposition}}  & A' \idPerp_G  B' \given C'\dcup\{v\} \\
      & \xRightarrow{\text{lemma}} &  A' \dcup \{v\} \idPerp_G  B' \given C' \qquad \lor \qquad A' \idPerp_G B' \dcup \{v\} \given C' \\
      & \xRightarrow{\text{cases (i),(ii)}} &  X_{A'},X_v \Indep_{R}  X_{B'} \given X_{C'} \quad \lor \quad X_{A'} \Indep_{R}  X_{B'},X_v \given X_{C'} \\
      & \xRightarrow{\text{weak union}} &  X_{A'} \Indep_{R}  X_{B'} \given X_{C} \quad \lor \quad X_{A'} \Indep_{R}  X_{B'} \given X_{C} \\
      & \xRightarrow{\text{redundancy}} &  X_A \Indep_{R}  X_B \given X_{C} \quad \lor \quad X_A \Indep_{R}  X_B \given X_{C}
    \end{align*}
    This shows the claim in case (iii).

  \item $v \notin A \cup B \cup C$. This follows from the induction hypothesis.
\end{enumerate}
\end{proof}
So the only properties of variation independence we needed were:
(i) stability under marginalization; (ii) local Markov property (a single childless node is independent of all others given its parents); (iii) (asymmetric) separoid axioms;
(iv) the global Markov property holds for $V=\emptyset$ and $J$ arbitrary. 

Beautiful, this gives a nice foundation for JCI with input nodes!


By the way: this allows us to strengthen this to $iD$-separation!

This is inspired by \cite{GVP1990}:
\begin{Def}
  A node $v$ is \emph{functionally determined} by $Z$ if $v \in Z$ or $v$ is a deterministic node and all its parents are functionally determined by $Z$.
  If $v$ is a deterministic node with no parents, then it is functionally determined by $Z$.
  A set of nodes is functionally determined by $Z$ if each of its members is determined by $Z$.
\end{Def}
Perhaps we can just say: $|R(X_v | X_Z)| = 1$.

\begin{Def}
  For $A,B,C$ disjoint, we say that $C$ \emph{D-separates} $A$ from $B$ if there exists no path between a node in $A$ and a node in $B$ along which (i) every collider is in $\Anc(C)$, (ii) every non-collider is not in $C$, and (iii) no tail-to-tail node on the path is functionally determined by $C$.
\end{Def}

Actually \cite{GVP1990} give quite a nice general exposition for separoids!
So before setting out to prove the corresponding Markov property yourself, see that paper!

Also, we may be able to derive stronger results because D-separation is not complete (see our \texttt{detsep.tex}).
